\section{Sources and legal basis}
\label{sec:legal}

\subsection{What was used}

Every statement in this paper derives from one of five kinds of source, and the
architecture in \S\ref{sec:arch} labels each parameter with the class of the
source it came from.

\begin{enumerate}
  \item \textbf{Open standards.} The InfiniBand Architecture Specification and
    its RoCEv2 annex~\cite{ibta-vol1,ibta-a17} define the transport, the base
    transport header and its extensions, the packet sequence number space, the
    acknowledgement and negative-acknowledgement semantics, the go-back-N
    recovery rule, and the encapsulation of an InfiniBand transport packet in
    UDP over IP over Ethernet. Everything the architecture does on the wire is
    \evdoc{} against these.
  \item \textbf{The public programmer's reference manual.} The vendor's
    published reference manual for this device family~\cite{mlx-prm} documents
    the software-visible objects: queue-pair context fields, work-queue-element
    layouts, completion-queue-entry formats, doorbell records, the user-access
    region, and the memory-translation and memory-protection tables. It
    documents an interface, not a microarchitecture, and we use it only as an
    interface document.
  \item \textbf{Open-source software.} The Linux \texttt{mlx5} RDMA and
    Ethernet drivers~\cite{linux-mlx5} and the \texttt{rdma-core}
    \texttt{mlx5} provider~\cite{rdma-core} are the authority for how work is
    actually submitted and completions actually read: linked work-request
    lists, work-queue-element construction, the publication barrier, the
    doorbell record update, the user-access-region write, inline and BlueFlame
    paths, completion-queue polling and completion decoding, and the names,
    widths and reset behaviour of the hardware counters. Behaviour that follows
    from reading this source but has no public silicon contract is \evdrv{}.
  \item \textbf{Published literature.} Peer-reviewed measurements and designs,
    cited where used, including the anomaly-search method we adapted, the
    congestion-control algorithm the device implements, and prior
    microarchitectural studies of RDMA NICs.
  \item \textbf{Black-box measurement.} More than 700 measured points on four
    cluster nodes that the authors are entitled to use, taken through the
    documented verbs interface and the documented counter interface, with no
    privileged access to the device: no firmware image was read, disassembled
    or modified, no non-public register was written, and no protection measure
    of any kind was circumvented. Behaviour whose aggregate effect is
    identifiable from these measurements but whose internal mechanism is not
    public is \evcal{}, and carries a named latent parameter with bounds rather
    than a fabricated register.
\end{enumerate}

No firmware was obtained, decompiled or inspected. No material under a
non-disclosure agreement was used, and none of the authors' NDA-covered
knowledge, where any exists, entered this document; the campaign records and
this paper were written from the five source classes above and nothing else.
The counters we read are the ones the open-source driver exposes through
\texttt{ethtool} and \texttt{sysfs} to any unprivileged user of the machine.
Several vendor configuration surfaces (the DCQCN parameter block, the firmware
non-volatile configuration, PCI capability space, kernel debug filesystem) were
found to be unreadable on the campaign hosts and are reported as unreadable
rather than guessed.

\subsection{Basis in United States law}

We record the authors' basis for the work, not legal advice, and we encourage
anyone reproducing it to take their own.

Clean-room reverse engineering of an interface and of externally observable
behaviour, from lawfully obtained public information and from a device the
investigator is entitled to use, for the purpose of interoperability and
research, is lawful in the United States. Two decisions of the Ninth Circuit
are directly on point. In \emph{Sega Enterprises Ltd.\ v.\ Accolade,
Inc.}~\cite{sega1992} the court held that disassembly of a lawfully acquired
program is fair use when it is the only means of gaining access to unprotected
functional elements and there is a legitimate reason for seeking that access.
In \emph{Sony Computer Entertainment, Inc.\ v.\ Connectix Corp.}~\cite{sony2000}
the court extended the reasoning to intermediate copying performed while
developing an independent, non-infringing implementation of a documented
interface. Congress codified a narrower but explicit permission in
17~U.S.C.~\S~1201(f)~\cite{usc1201f}, which allows circumvention and the
development of the means to circumvent for the sole purpose of achieving
interoperability of an independently created computer program, to the extent
that the acts do not constitute infringement.

Our position is more conservative than any of these authorities requires. We
did not disassemble anything, we made no intermediate copy of a protected work,
and we circumvented no technological protection measure, so \S~1201(f) is
cited as the boundary of what would have been permissible rather than as the
permission we relied on. What we did is read public documents, read
open-source code under its own licence, run standard user-space benchmarks and
a modified open-source traffic generator against hardware the authors are
entitled to use, and read counters that the operating system exposes. The
functional facts recovered by that process (that a device stamps ECT(0), that a
counter is inert, that an ingress meter drains at a given rate) are facts, and
facts are not copyrightable. The architecture in \S\ref{sec:arch} is an
independent design that reproduces measured behaviour; it is not a copy of any
vendor design, and where the mechanism is unknown it says so rather than
inventing a plausible one.

\subsection{What we publish}

The campaign records, the curated measurement extracts plotted here, the golden
C model, the anomaly table and the architecture are all public. The raw bulk
captures (per-packet logs on the order of $10^8$ packets) stay outside the
repository for size reasons and are available on request. Node names and
addresses are removed from every published extract.
